\documentclass{article}

\usepackage{amsmath}
\usepackage{amsthm}
\theoremstyle{definition}
\newtheorem{definition}{Definition}
\newtheorem{lemma}[definition]{Lemma}

\title{Ideas for a semantics for HiFIRRTL}
\author{David N. Jansen \\ based on discussion with Xiaomu Shi, Keyin Wang and Zhilin Wu}

\begin{document}

\maketitle

\section{Introduction}

This document describes a semantics for HiFIRRTL modules and circuits.
The basic idea is that every HiFIRRTL module is mapped to a \emph{module graph.}
Actually this is a hypergraph, because edges may have a tree structure as well.
Module graphs can be combined into a circuit graph,
similar to how HiFIRRTL modules are combined into a circuit.

The transformation from HiFIRRTL to LoFIRRTL should preserve the semantics on the level of module graphs
as much as possible.
Only for a few passes, we need to look into the behavioural semantics of a module graph,
namely when the structure is replaced by a bisimilar structure.

\section{Module graphs}

\begin{definition}
A module graph consists of:
\begin{itemize}
\item	a set of components.
	There are the following components:
	input ports;
	output ports;
	wires;
	registers;
	(perhaps nodes, but I don't think that is necessary);
	memories;
	instances;
	constants;
	(perhaps invalid sources, to translate an ``is invalid'' statement);
	multiplexers;
	primitive operations (add, subtract, multiply, divide, modulus, comparison, padding,
		type cast, shift left, shift right, dynamic shift left, dynamic shift right,
		convert to SInt, negate, bitwise complement, bitwise and, bitwise or, bitwise xor,
		reduce by and, reduce by or, reduce by xor, concatenate, extract bits);
	an instance of another module graph.

	Every component has some input connectors and/or output connectors.
	Input and output connectors have a data type (e.g.\@ $\mathsf{SInt}$) and a bit width.
\item	a set of connection trees.
	A connection tree links any number of output connectors of component(s), called sources,
	to an input connector, called destination.
	If multiple sources are used, then a ternary tree structure indicates
	how to decide which one to use in a specific condition:
	every internal node of the tree has structure $(\mathit{cond}, \mathit{connect}_\mathit{true}, \mathit{connect}_\mathit{false})$.
	Depending on $\mathit{cond}$, one of the two connectors is actually used.
	Note that $\mathit{cond}$ is also linked to an output connector of a component.
	A connection tree has a data type (e.g. $\mathsf{SInt}$) and a bit width.
\end{itemize}
\end{definition}

\noindent
A module graph has a behavioural semantics,
where inputs and outputs behave as in (idealized) circuits
and connection trees decide dynamically based on the respective condition
which connector to use.

\section{Semantics of HiFIRRTL modules}

The structural semantics of a HiFIRRTL module is defined to be a set of module graphs.

Given a module with input ports $I$, output ports $O$, and statement sequence $S$,
the respective module graph has to contain:
\begin{itemize}
\item	an input port component for every element in $I$.
\item	an output port component for every element in $O$.
\item	other components and connection trees are defined inductively depending on $S$:
	\begin{itemize}
	\item	If $S$ is the empty statement sequence,
		then there are no (further) components or connection trees.
	\item	If $S = (s_1, s_2, \ldots, s_{n+1})$,
		then the module graph is a modification of the graph $G$ for $I, O, (s_1, \ldots, s_n)$ as follows:
		\begin{itemize}
		\item	If $s_{n+1} = \mathtt{skip}$, the result is $G$.
		\item	If $s_{n+1} = \mathtt{register} \ldots$, the result is $G$ with an additional register component as specified in $s_{n+1}$.
		\item	If $s_{n+1} = \mathit{target} \mathtt{<=} \mathit{source}$, the result is $G$ with an additional or changed connection tree to $\mathit{target}$,
			and possibly components as indicated in the expression $\mathit{source}$.
			(If defined in detail, this will require an inductive semantics of expressions.)
		\item	similarly for other statements
		\end{itemize}
	\end{itemize}
	The bit widths of every input and output connector satisfy the following conditions:
	If the HiFIRRTL module specifies a bit width explicitly,
	then that is the bit width of the connector.
	Otherwise, any bit width is allowed,
	as long as it satisfies the default bit width constraints of all components in the HiFIRRTL module.
\end{itemize}

\noindent
If all bit widths in the HiFIRRTL module are specified,
then there is at most one module graph that can be the semantics of the HiFIRRTL module;
otherwise, there may be multiple module graphs,
namely all those that satisfy the default bit width constraints.
(An ill-specified HiFIRRTL module may have no semantics at all,
namely if the bit widths of some connection trees cannot be determined at all.)

What are the default bit width constraints?
Every HiFIRRTL component specifies a default width for its output connectors,
which depends on the input connector widths.
For example, for an adder, the specification is
$\mathit{width}_\mathit{result} \geq \max \{\mathit{width}_a, \mathit{width}_b\} + 1$.

Note that a module graph does not contain aggregate type components or connections.
So the HiFIRRTL semantics also has to translate any component with an aggregate type
into multiple components in the graph.
\textbf{TO DO:} This may be problematic for subaccesses, dynamic accesses to an array.
Perhaps we have to allow arrays after all.

\section{Remarks on the lowering transforms}

\subsection{InferWidths}

chooses the minimal element (in terms of bit width) from the set of module graphs and changes the HiFIRRTL module such that only this element remains in the semantics.

\begin{definition}
	The bit-width-preorder between module graphs $G_1, G_2$ is defined as follows:
	$G_1 \leq G_2$ if the two module graphs contain the same elements,
	and the bit width of each element in $G_1$ is smaller than or equal to the bit width of the same element in $G_2$.
\end{definition}

\noindent
Based on a lemma like the one below
we can guarantee that the minimal element in the set of module graphs does not incur any unintended data loss.
(Data loss may occur if the user chooses a non-default bit width,
e.g.\@ if the result of an addition of 32-bit summands is stored in a 32-bit register,
but then the user has to explicitly specify that the goal register has 32 bits.)

\begin{lemma}
	The output bit widths of every HiFIRRTL component are monotonous,
	i.e.\@ if no data loss occurs at output bit width $n$,
	then no data loss occurs at any output bit width $>n$.
\end{lemma}

\subsection{LowerTypes}

changes the HiFIRRTL module code but not its semantics
because aggregate-type components were already translated to sets of ground-type components before.

\subsection{ExpandWhens}

replaces the connection trees by multiplexers.
We will need to prove that this does not fundamentally change the behaviour
through a bisimulation between the module graphs.

\end{document}
